\begin{titlepage}
	\definecolor{forestgreen}{RGB}{34,92,56}   % 森林绿主题色
	\begin{tikzpicture}[remember picture, overlay]
		% ---- 1. 墨绿渐变铺满整个封面 ----
		\shade[top color=forestgreen, middle color=forestgreen!90, bottom color=forestgreen!60!black]
			(current page.south west) rectangle (current page.north east);

		% ---- 2. 细点阵（模拟网格，微弱纹理）----
		\foreach \x in {1,...,20} {
			\foreach \y in {1,...,29} {
				\fill[white, opacity=0.05] (\x,\y) circle[radius=0.018cm];
			}
		}

		% ---- 3. 主标题后方的柔光晕 ----
		\fill[white, opacity=0.05] ([yshift=-8.2cm]current page.north) circle[radius=7.0cm];
		\fill[white, opacity=0.05] ([yshift=-8.2cm]current page.north) circle[radius=4.6cm];

		% ---- 4. 右上：复平面单位圆（复数 z = e^{iθ}）----
		\begin{scope}[shift={([xshift=-6.2cm,yshift=-4.7cm]current page.north east)}]
			% 复平面坐标轴
			\draw[white, opacity=0.6, line width=0.8pt, -{Stealth}] (0.5,1.5) -- (4.7,1.5);
			\draw[white, opacity=0.6, line width=0.8pt, -{Stealth}] (2.6,0.3) -- (2.6,3.0);
			\node[anchor=north, text=white, opacity=0.8, font=\footnotesize\itshape] at (4.7,1.5) {$x$};
			\node[anchor=east, text=white, opacity=0.8, font=\footnotesize\itshape] at (2.6,3.05) {$y$};
			\node[anchor=north east, text=white, opacity=0.7, font=\footnotesize\itshape] at (2.6,1.42) {$O$};
			% 单位圆
			\draw[white, line width=1.1pt, opacity=0.9] (2.6,1.5) circle (1.15cm);
			% 向量到点 z
			\draw[white, line width=1.0pt, opacity=0.95, -{Stealth}] (2.6,1.5) -- (3.30,2.38);
			\fill[cyan!65!white, opacity=0.95] (3.30,2.38) circle[radius=0.05cm];
			% 角度弧
			\draw[white, dashed, opacity=0.5, line width=0.5pt] (2.9,1.5) arc (0:52:0.3);
			\node[anchor=south west, text=white, opacity=0.9, font=\footnotesize\itshape] at (2.75,2.55) {$z=e^{i\theta}$};
		\end{scope}

		% ---- 5. 左下：指数衰减正弦曲线（Fourier / Laplace 主题）----
		\begin{scope}[shift={([xshift=0.9cm,yshift=0.9cm]current page.south west)}]
			% 坐标轴
			\draw[white, opacity=0.6, line width=0.8pt, -{Stealth}] (0.2,1.4) -- (3.4,1.4);
			\draw[white, opacity=0.6, line width=0.8pt, -{Stealth}] (0.2,0.3) -- (0.2,3.0);
			% 包络虚线
			\draw[white, dashed, opacity=0.35, line width=0.6pt, domain=0.1:3.2, samples=60, smooth]
				plot (\x, {1.4 + 1.1*exp(-0.3*\x)});
			\draw[white, dashed, opacity=0.35, line width=0.6pt, domain=0.1:3.2, samples=60, smooth]
				plot (\x, {1.4 - 1.1*exp(-0.3*\x)});
			% 指数衰减正弦
			\draw[white, line width=1.2pt, opacity=0.95, domain=0.1:3.2, samples=120, smooth]
				plot (\x, {1.4 + 1.1*exp(-0.3*\x)*sin(deg(3.5*\x))});
		\end{scope}

		% ---- 6. 金色细边框 ----
		\draw[gold, line width=0.7pt, opacity=0.55]
			([shift={(0.85,0.85)}]current page.south west) rectangle
			([shift={(-0.85,-0.85)}]current page.north east);

		% ---- 7. 顶部眉题（英文小字 + 金色短线）----
		\coordinate (T) at ([yshift=-3.5cm]current page.north);
		\draw[gold, line width=1.0pt] (T) ++(-3.4cm,0) -- ++(1.0cm,0);
		\draw[gold, line width=1.0pt] (T) ++(2.4cm,0) -- ++(1.0cm,0);
		\node[anchor=north, text=white, opacity=0.85, align=center] at (T) {
			{\sffamily\fontsize{13}{16}\selectfont
				\uppercase{Engineering \ Mathematics}}
		};

		% ---- 8. 主标题（居中，使用 \notename）----
		\node[anchor=center, align=center, text=white] at ([yshift=-8.0cm]current page.north) {%
			{\fontsize{40}{50}\sffamily\bfseries \notename}
		};

		% ---- 9. 金色双细线 + 单位圆纹饰 ----
		\coordinate (C) at ([yshift=-11.4cm]current page.north);
		\draw[gold, line width=1.1pt] (C) ++(-2.2cm,0) -- ++(1.55cm,0);
		\draw[gold, line width=1.1pt] (C) ++(0.65cm,0) -- ++(1.55cm,0);
		\draw[gold, line width=0.9pt] (C) circle[radius=0.16cm];
		\draw[gold, line width=0.9pt, opacity=0.85] (C) circle[radius=0.05cm];

		% ---- 10. 副标题 ----
		\node[anchor=north, text=white, opacity=0.9, align=center]
			at ([yshift=-12.2cm]current page.north) {
			{\fontsize{16}{20}\sffamily 学\ \ 习\ \ 笔\ \ 记}
		};

		% ---- 11. 底部作者、日期信息 ----
		\coordinate (B) at ([yshift=3.4cm]current page.south);
		\draw[gold, line width=0.8pt, opacity=0.7] (B) ++(-1.1cm,0.85cm) -- ++(2.2cm,0);
		\node[anchor=south, text=white, align=center] at (B) {
			\begin{tabular}{c}
				{\fontsize{20}{24}\sffamily\bfseries \noteauthor} \\[0.4cm]
				{\large \sffamily 2026年6月10日}
			\end{tabular}
		};
	\end{tikzpicture}
\end{titlepage}
